\chapter*{Introduction générale}
\addcontentsline{toc}{chapter}{Introduction générale}

La transformation numérique constitue aujourd'hui un levier important pour améliorer la gestion et l'efficacité des entreprises. Dans le secteur automobile, la diversité des activités, allant de la commercialisation des pièces à la réalisation des services de maintenance, nécessite des outils permettant de centraliser les informations et de faciliter le suivi des opérations. C'est dans ce contexte que ce projet a été réalisé pour \textbf{MSG}, un garage spécialisé dans la vente de pièces automobiles et les services de maintenance, pour objectif de lui proposer une solution digitale adaptée à ses besoins et destinée à être intégrée à son activité.

La gestion simultanée des produits, des stocks, des fournisseurs, des commandes, des clients et des interventions de maintenance peut devenir complexe lorsque les informations sont réparties entre plusieurs outils ou traitées manuellement. Cette situation peut entraîner des difficultés dans le suivi des opérations, la circulation des informations et la coordination entre les différents intervenants. Il devient donc nécessaire de disposer d'une plateforme centralisée permettant d'unifier les principaux processus de l'entreprise et de faciliter leur gestion, notamment la gestion des utilisateurs, des rôles et des autorisations, la gestion des produits, des stocks et des fournisseurs, la gestion des commandes, des factures et des paiements, la gestion des services de maintenance et des interventions, la gestion des employés et de certaines opérations liées aux ressources humaines, l'automatisation de certaines tâches et notifications, ainsi qu'une interface e-commerce destinée aux clients.

Dans ce contexte, se situe notre solution par le biais de ce projet qui vise à concevoir et développer \textbf{AutoParts Pro}, une plateforme combinant les fonctionnalités d'un \textbf{site e-commerce} et d'un \ac{ERP} destiné à la gestion interne de \textbf{MSG}. La solution permet de centraliser les processus liés à la vente de pièces automobiles, à la gestion des stocks et des fournisseurs, aux commandes et aux paiements, ainsi qu'aux services de maintenance, tout en intégrant un mécanisme de gestion des utilisateurs et des autorisations afin de sécuriser l'accès aux différentes fonctionnalités. L'objectif est ainsi de fournir à \textbf{MSG} une solution évolutive permettant d'améliorer l'organisation de ses activités, de centraliser les informations et d'offrir aux clients une expérience numérique adaptée.

Le présent rapport décrit en détail la progression du projet qui s'étale sur sept chapitres comme suit : le premier chapitre présente l'organisme d'accueil, le client \textbf{MSG}, la motivation et la problématique du projet, l'étude de l'existant ainsi que la méthodologie Scrum adoptée et la planification du travail. Le deuxième chapitre présente la spécification des besoins, à travers l'identification des acteurs et des exigences fonctionnelles et non fonctionnelles, les environnements techniques utilisés ainsi que l'architecture générale de la plateforme. Les chapitres trois à sept sont consacrés à la réalisation des différents Sprints, couvrant respectivement la gestion des utilisateurs et de l'authentification, la gestion des produits et du stock, les services de maintenance, la gestion des ressources humaines, ainsi que l'intégration et l'automatisation avec n8n. Enfin, nous clôturons notre rapport par une conclusion générale exposant une synthèse du travail ainsi que les principaux résultats obtenus, les limites éventuelles de la solution et quelques perspectives d'évolution.